% GENERATED FILE. Edit canonical lessons, then run npm run generate:books.
% canonical-source-hash: fnv1a64:7b73a6da1cb63373
% canonical-lessons: ZH-C406-yongji, ZH-C406-laihui, ZH-C406-gaiqian, ZH-C406-qiche, ZH-C406-xinglijia

\chapter{A Busy Platform}
\label{ch:zh-a-busy-platform}
\hlchaptermodality{406}

\begin{chapteropening}
\textbf{By the end of this chapter:} \emph{I can say a place is crowded, talk about a round trip, change a booking and talk about cars and racks.}

Complete the last lesson of chapter 406: I can say a place is crowded, talk about a round trip, change a booking and talk about cars and racks.
\end{chapteropening}

\section[yōngjǐ]{yōngjǐ — crowded}

\label{lesson:ZH-C406-yongji}



\begin{quote}
Before the new one: say the Mandarin for the last bus or train, then the Mandarin for to call at a station.
\end{quote}

\subsection*{You'll want to know: yōngjǐ}
\emph{yōngjǐ} — \textquotedblleft{}crowded\textquotedblright{}.

\begin{quote}
\emph{dìtiě tài yōngjǐ le wǒmen zuò chūzūchē qù} — The metro's too crowded: we'll take a taxi.
\end{quote}

Its characters come later; for now it is a word for the ear.

\begin{cousinweb}
\emph{yōngjǐ} is built from \emph{yōng} 'to hold close' and \emph{jǐ} 'to squeeze', both inherited from Middle Chinese, \emph{'jowngX} and \emph{tsej} (Guangyun, in Baxter's transcription), which glosses the first 'to hold in the arms' and the second 'to push aside'. Those forms would regularly give yǒng and jī: the tones of yōngjǐ are irregular.
\end{cousinweb}

\begin{grammarlens}[title={What you've built}]
One more word for getting around.
\end{grammarlens}

\subsection*{Your turn}
\begin{itemize}
\raggedright
  \item \emph{Say it:} \emph{yōngjǐ}
  \item \emph{Say it:} \emph{yōngjǐ}, once more
  \item \emph{Say it:} \emph{dìtiě tài yōngjǐ le wǒmen zuò chūzūchē qù} — The metro's too crowded: we'll take a taxi.
  \item \emph{Recall:} say the Mandarin for the last bus or train, then the Mandarin for to call at a station, then say \emph{yōngjǐ} again
\end{itemize}

\begin{tcolorbox}[breakable,colback=teal!4,colframe=teal!35!black,title={Before you move on}]
What does yōngjǐ mean? (\textquotedblleft{}Crowded\textquotedblright{}.) Say it once more.
\end{tcolorbox}
\section[\texorpdfstring{\zh{来回} (láihuí)}{来回 (láihuí)}]{\zh{来回} (láihuí) — round trip}

\label{lesson:ZH-C406-laihui}



\begin{quote}
Before the new one: say the Mandarin for to call at a station, then the Mandarin for crowded.
\end{quote}

\subsection*{You'll want to know: \zh{来回}}
\textbf{\zh{来回}} — \emph{láihuí} — \textquotedblleft{}round trip\textquotedblright{}.

\begin{quote}
\emph{cóng zhèr dào jīchǎng yí ge láihuí yào liǎng ge xiǎoshí} — A round trip from here to the airport takes two hours.
\end{quote}

\begin{cousinweb}
\textbf{\zh{来回}} (\emph{láihuí}) is built from \zh{来} 'to come' and \zh{回} 'to go back', both inherited from Middle Chinese, \emph{loj} and \emph{hwoj} (Guangyun, in Baxter's transcription), which glosses the first 'to arrive' and the second 'to turn'.
\end{cousinweb}

\begin{grammarlens}[title={What you've built}]
One more word for getting around.
\end{grammarlens}

\subsection*{Your turn}
\begin{itemize}
\raggedright
  \item \emph{Say it:} \emph{láihuí}
  \item \emph{Say it:} \emph{láihuí}, once more
  \item \emph{Say it:} \emph{cóng zhèr dào jīchǎng yí ge láihuí yào liǎng ge xiǎoshí} — A round trip from here to the airport takes two hours.
  \item \emph{Recall:} say the Mandarin for to call at a station, then the Mandarin for crowded, then say \emph{láihuí} again
\end{itemize}

\begin{tcolorbox}[breakable,colback=teal!4,colframe=teal!35!black,title={Before you move on}]
What does \zh{来回} mean? (\textquotedblleft{}Round trip\textquotedblright{}.) Say it once more.
\end{tcolorbox}
\section[gǎiqiān]{gǎiqiān — to change a booking; to change a ticket to another train}

\label{lesson:ZH-C406-gaiqian}



\begin{quote}
Before the new one: say the Mandarin for crowded, then the Mandarin for a round trip.
\end{quote}

\subsection*{You'll want to know: gǎiqiān}
\emph{gǎiqiān} — \textquotedblleft{}to change a booking; to change a ticket to another train\textquotedblright{}.

\begin{quote}
\emph{wǒ xiǎng gǎiqiān míngtiān de huǒchē} — I'd like to change to tomorrow's train.
\end{quote}

Its characters come later; for now it is a word for the ear.

\begin{cousinweb}
\emph{gǎiqiān} is built from \emph{gǎi} 'to alter' and \emph{qiān} 'to sign'. The first is inherited from Middle Chinese \emph{kojX} (Guangyun, in Baxter's transcription), glossed 'to change'; no dictionary we use records an older form of the second, so its history is not settled.
\end{cousinweb}

\begin{grammarlens}[title={What you've built}]
One more word for getting around.
\end{grammarlens}

\subsection*{Your turn}
\begin{itemize}
\raggedright
  \item \emph{Say it:} \emph{gǎiqiān}
  \item \emph{Say it:} \emph{gǎiqiān}, once more
  \item \emph{Say it:} \emph{wǒ xiǎng gǎiqiān míngtiān de huǒchē} — I'd like to change to tomorrow's train.
  \item \emph{Recall:} say the Mandarin for crowded, then the Mandarin for a round trip, then say \emph{gǎiqiān} again
\end{itemize}

\begin{tcolorbox}[breakable,colback=teal!4,colframe=teal!35!black,title={Before you move on}]
What does gǎiqiān mean? (\textquotedblleft{}To change a booking; to change a ticket to another train\textquotedblright{}.) Say it once more.
\end{tcolorbox}
\section[qìchē]{qìchē — motor vehicle; car}

\label{lesson:ZH-C406-qiche}



\begin{quote}
Before the new one: say the Mandarin for a round trip, then the Mandarin for to change a booking.
\end{quote}

\subsection*{You'll want to know: qìchē}
\emph{qìchē} — \textquotedblleft{}motor vehicle; car\textquotedblright{}.

\begin{quote}
\emph{wǒmen zuò qìchē qù jīchǎng} — We're going to the airport by car.
\end{quote}

Its characters come later; for now it is a word for the ear.

\begin{cousinweb}
\emph{qìchē} is built from \emph{qì} 'steam' and \zh{车} 'vehicle'. The second is inherited from Middle Chinese \emph{tsyhae} (Guangyun, in Baxter's transcription); no dictionary we use records an older form of the first, so its history is not settled.
\end{cousinweb}

\begin{grammarlens}[title={What you've built}]
One more word for getting around.
\end{grammarlens}

\subsection*{Your turn}
\begin{itemize}
\raggedright
  \item \emph{Say it:} \emph{qìchē}
  \item \emph{Say it:} \emph{qìchē}, once more
  \item \emph{Say it:} \emph{wǒmen zuò qìchē qù jīchǎng} — We're going to the airport by car.
  \item \emph{Recall:} say the Mandarin for a round trip, then the Mandarin for to change a booking, then say \emph{qìchē} again
\end{itemize}

\begin{tcolorbox}[breakable,colback=teal!4,colframe=teal!35!black,title={Before you move on}]
What does qìchē mean? (\textquotedblleft{}Motor vehicle; car\textquotedblright{}.) Say it once more.
\end{tcolorbox}
\section[xínglijià]{xínglijià — luggage rack}

\label{lesson:ZH-C406-xinglijia}



\begin{quote}
Before the new one: say the Mandarin for to change a booking, then the Mandarin for a motor car.
\end{quote}

\subsection*{You'll want to know: xínglijià}
\emph{xínglijià} — \textquotedblleft{}luggage rack\textquotedblright{}.

\begin{quote}
\emph{zhège xínglijià hěn gāo} — This luggage rack is very high.
\end{quote}

Its characters come later; for now it is a word for the ear.

\begin{cousinweb}
\emph{xínglijià} is built from \emph{xíngli} 'luggage' and \emph{jià} 'a frame', all inherited from Middle Chinese, \emph{haeng}, \emph{liX} and \emph{kaeH} (Guangyun, in Baxter's transcription), which glosses the last as the frame of a house.
\end{cousinweb}

\begin{grammarlens}[title={What you've built}]
One more word for getting around.
\end{grammarlens}

\subsection*{Your turn}
\begin{itemize}
\raggedright
  \item \emph{Say it:} \emph{xínglijià}
  \item \emph{Say it:} \emph{xínglijià}, once more
  \item \emph{Say it:} \emph{zhège xínglijià hěn gāo} — This luggage rack is very high.
  \item \emph{Recall:} say the Mandarin for to change a booking, then the Mandarin for a motor car, then say \emph{xínglijià} again
\end{itemize}

\begin{tcolorbox}[breakable,colback=teal!4,colframe=teal!35!black,title={Before you move on}]
What does xínglijià mean? (\textquotedblleft{}Luggage rack\textquotedblright{}.) Say it once more.
\end{tcolorbox}
